% ---------------------------------------------------------------------
%  appendix-a-formal.tex -- the formal HPB model
%  Single place where Definitions, Rules R1--R7 (labels rule:*) and
%  Properties P1--P4 (labels prop:*) are stated. The main text
%  (sections 03, 04, 05) explains them in plain language.
% ---------------------------------------------------------------------
\section{Formal Model}
\label{app:formal}

This appendix states HPB precisely, for readers who want to check the
method or build a tool that enforces it. \Cref{sec:model} explains the same
concepts in plain language; the rules R1--R7 are introduced, with their
practical motivation, in \cref{sec:onepage}. The last three subsections
give the exact definitions of the metrics, the layout algorithm, and the
data model and status of the prototype. Time is taken from a totally
ordered domain $\mathbb{T}$ (calendar time); how time is \emph{drawn} is a
separate matter (\cref{sec:notation,def:ordinal}). Concepts that only some
organizations need (configurations, dependencies, launch bundles and the
deployment overlay) can be left out without affecting the rest of the
model.

\subsection{Lines and baselines}
\label{app:formal-lines}

\begin{definition}[Product line, line forest]
  \label{def:line}
  A \emph{product line} $P$ is an identified product reference, such as a
  model together with its manufacturer reference, or a firmware or
  software product. A portfolio has a finite set of lines $\mathcal{P}$
  and a partial function
  $\mathrm{parent}\colon \mathcal{P} \rightharpoonup \mathcal{P}$, where
  $\mathrm{parent}(P') = P$ states that $P'$ was created by a hardware
  revision of $P$ (\cref{def:linefork}). The graph with an edge from
  $\mathrm{parent}(P')$ to $P'$ is a forest, the \emph{line forest}.
\end{definition}

\begin{definition}[Baseline]
  \label{def:baseline}
  A \emph{baseline} $v$ is a named, identifiable state of a product line
  $\ell(v) \in \mathcal{P}$. It carries a name, unique within its line
  (e.g., \texttt{v2.3.1} or \texttt{Rev~B}), a kind $k(v)$, one of
  \textsf{firmware}, \textsf{software}, \textsf{hardware} or
  \textsf{other}, a status $s(v)$, either \textsf{planned} or
  \textsf{released}, and a date $t(v) \in \mathbb{T}$, which is the
  release date if $v$ is released and the target date if it is planned.
  We write $\mathcal{V}$ for the set of all baselines and
  $\mathcal{V}(P) = \{v \in \mathcal{V} : \ell(v) = P\}$. A status may
  change only from \textsf{planned} to \textsf{released}; once released,
  a baseline is immutable.
\end{definition}

The kind \textsf{other} is used for process baselines. The normative
release event per kind, and the form--fit--function criterion that decides
whether a hardware change is a baseline of an existing line or a line
fork, are given in \cref{sec:model-lines}.

\begin{definition}[Release stream, main line, precedence]
  \label{def:mainline}
  The baselines of a line $P$ are partitioned into one or more
  \emph{release streams}. The \emph{main stream} starts at the
  \emph{first baseline} $v_0(P)$ of $P$, i.e., its first release or, for a
  line created by a line fork, the hardware baseline that created it
  (\cref{def:linefork}). Every other stream $R$, a \emph{maintenance
  stream}, is created at a released \emph{stream-fork baseline}
  $\beta(R)$ of another stream of $P$. Within a stream, released baselines
  are ordered by date and followed by the planned baselines of the stream,
  ordered by target date; the sequence of the main stream,
  \[
    M(P) = (v_0, v_1, \dots, v_n, u_1, \dots, u_r),
  \]
  is the \emph{main line} of $P$. \emph{Precedence} $\prec$ is the smallest
  transitive relation on $\mathcal{V}$ such that $x \prec y$ whenever
  (i)~$y$ immediately follows $x$ in a stream, (ii)~$y$ is the first
  baseline of a maintenance stream $R$ and $x = \beta(R)$, or (iii)~$y$ is
  the first baseline of a child line $P'$ and $x = \delta(P')$, the
  baseline of the parent line from which $P'$ was derived
  (\cref{def:linefork}). We write $\preceq$ for its reflexive closure.
\end{definition}

By construction every baseline has at most one immediate predecessor: its
predecessor in its stream, the fork baseline of its stream, or the parent
baseline of its line. Precedence is therefore a forest order, and in
particular a strict partial order. Release dates order baselines only
within a stream.

\begin{definition}[Applies-to relation, configuration]
  \label{def:config}
  A line is a \emph{hardware} line if its reference is a physical unit and
  a \emph{software} line if it is a firmware, application or service
  product; a \emph{combined} line carries the hardware revisions of one
  reference together with the firmware released for it. For every
  baseline $v$ of a software line, the \emph{applies-to} set
  $\mathrm{app}(v) \subseteq \mathcal{P}$ lists the hardware references on
  which $v$ is released to run (its effectivity). A \emph{configuration}
  $\kappa$ is a set of baselines with at most one baseline per line in
  which every software baseline applies to the hardware reference of
  $\kappa$; for example (\texttt{Rev~A}, firmware \texttt{v2.0}, app
  \texttt{5.1}). The relation $\mathrm{app}$ is the compatibility matrix
  of the portfolio.
\end{definition}

\subsection{Branches, states and closure}
\label{app:formal-branches}

\begin{definition}[Branch]
  \label{def:branch}
  A \emph{branch} $b$ is a unit of change intent on a product. We write
  $\mathcal{B}$ for the set of all branches. Every branch has
  \begin{itemize}[leftmargin=1.5em,itemsep=1pt,topsep=2pt]
    \item a \emph{type} $\tau(b) \in \{\mathsf{I}, \mathsf{X}, \mathsf{T},
          \mathsf{L}\}$: improvement, issue (defect), task, or laboratory
          evolution;
    \item an \emph{opening time} $t_o(b) \in \mathbb{T}$, the time at which
          the request is accepted at triage;
    \item a \emph{fork baseline} $f(b) \in \mathcal{V}$ with
          $s(f(b)) = \textsf{released}$ and $t(f(b)) \le t_o(b)$, fixed at
          intake: for an issue, the baseline in which the defect was
          \emph{observed}, as reported; for other types, the baseline
          relative to which the change is specified, by default the latest
          released baseline of the main stream at $t_o(b)$;
    \item for an issue, optionally, the \emph{observed configuration}
          $\Phi(b) \ni f(b)$ when the defect was reported for a combination
          of baselines of several lines; $f(b)$ is then its \emph{primary}
          element, on whose line the branch is drawn and attributed;
    \item for an issue, two attributes set by later analysis: the set
          $\mathrm{aff}(b) \subseteq \mathcal{V}$ of further baselines in
          which the defect is confirmed (\emph{also affects}) and, when root
          cause is known, the baseline $\iota(b)$ that \emph{introduced} it;
    \item a finite, time-ordered sequence of \emph{steps}
          $S(b) = \big((t_1, \kappa_1, e_1), \dots, (t_q, \kappa_q, e_q)\big)$
          with $t_o(b) \le t_1$, kinds
          $\kappa_i \in \{\textsf{done}, \textsf{todo}, \textsf{pass},
          \textsf{fail}\}$ and optional evidence $e_i$ (a report, a
          measurement, the identifier of a build).
  \end{itemize}
\end{definition}

The fork baseline is never revised by later knowledge: confirmation in
other baselines extends $\mathrm{aff}(b)$ and root-cause analysis sets
$\iota(b)$. An error in $f(b)$ itself is corrected only through an audited
transition, and an inferred fork baseline is flagged as such.

\begin{definition}[Branch state, verification]
  \label{def:state}
  The \emph{state} of a branch is a function
  $\sigma\colon \mathcal{B} \times \mathbb{T} \to \{\textsf{open},
  \textsf{targeted}, \textsf{merged}, \textsf{discarded}\}$, defined for
  $t \ge t_o(b)$. A \emph{targeted} branch is open and committed to a
  \emph{target baseline} $g(b)$, which must be planned when the target is
  set. States change only through the transitions of \cref{fig:lifecycle}:
  \emph{target} and \emph{re-target} (set or change $g(b)$),
  \emph{un-target} (clear $g(b)$), \emph{merge}, \emph{discard} and, as an
  exceptional and audited operation, \emph{reopen}. The states
  \textsf{merged} and \textsf{discarded} are \emph{terminal}; the time of
  the most recent terminal transition is the \emph{closing time} $t_c(b)$,
  and $t_c(b) = \infty$ while $b$ is open or targeted. A targeted branch is
  \emph{verified} at time $t$ if $S(b)$ contains a \textsf{pass} step at a
  time $\le t$ whose evidence identifies a candidate of $g(b)$ (a release
  candidate build, a prototype or a pilot lot); the time of the latest such
  step is its \emph{verification time} $t_r(b)$.
\end{definition}

\begin{definition}[Merge and discard]
  \label{def:merge}
  \emph{Merging} a targeted branch $b$ sets its \emph{merge baseline}
  $m(b) = g(b)$ and establishes the relation $\mathrm{closes}(b, m(b))$:
  baseline $m(b)$ contains the resolution of $b$. Merges happen only when
  $m(b)$ is published (\cref{def:gate}), so $t_c(b) = t(m(b))$ for every
  merged branch. \emph{Discarding} a branch closes it without a baseline
  and records a non-empty \emph{rationale} $\rho(b)$ from one of two
  classes: \emph{abandoned} (not viable, duplicate of $b'$, cannot
  reproduce, superseded, out of scope, will not fix) or \emph{resolved
  without a product baseline} (already resolved in $v$, contained without
  design change, resolved by the supplier).
\end{definition}

\subsection{Line forks, life cycle, gates and relations}
\label{app:formal-forks}

\begin{definition}[Line fork]
  \label{def:linefork}
  A \emph{line fork} is a baseline $h$ with $k(h) = \textsf{hardware}$ that
  creates a new product reference $P'$, with $\mathrm{parent}(P') = P$.
  The baseline $h$ belongs to the child line, $\ell(h) = P'$, and is its
  first baseline, $v_0(P') = h$. The \emph{derived-from baseline}
  $\delta(P') \in \mathcal{V}(P)$ is the latest released baseline of the
  main stream of $P$ before $t(h)$; it is the only immediate predecessor of
  $h$ (\cref{def:mainline}).
\end{definition}

Branches that produce the new reference fork from baselines of $P$ and
merge into $h$, which is admitted by the precedence edge
$\delta(P') \prec h$. Because $h$ is not a baseline of the parent line,
later baselines of $P$ do not contain that work.

\begin{definition}[Line life cycle]
  \label{def:lifecycle}
  Every line $P$ has a sequence of dated \emph{milestones}
  \[
    t_{\mathrm{sale}}(P) \le t_{\mathrm{prod}}(P) \le t_{\mathrm{sec}}(P)
    \le t_{\mathrm{eos}}(P)
  \]
  (end of sale, end of production, start of security-only maintenance,
  end of service), each possibly undefined while it has not been decided.
  They divide the life of the line into the phases \emph{active},
  \emph{no longer sold}, \emph{no longer produced}, \emph{security only}
  and \emph{retired}. \Cref{tab:phases} states which branches may fork from
  the baselines of a line, and which baselines it may receive, in each
  phase.
\end{definition}

\begin{table}[tbp]
  \centering
  \small
  \caption{What a line still accepts in each phase of its life
    (\cref{def:lifecycle}). These are defaults: an organization may
    tighten them; relaxing them is an explicit decision recorded on the
    line.}
  \label{tab:phases}
  \begin{tabularx}{\linewidth}{@{}>{\raggedright\arraybackslash}p{3.2cm}
      >{\raggedright\arraybackslash}X
      >{\raggedright\arraybackslash}X@{}}
    \toprule
    \textbf{Phase (from milestone)} & \textbf{Branches that may fork}
      & \textbf{Baselines that may be added} \\
    \midrule
    Active & $\mathsf{I}$, $\mathsf{L}$, $\mathsf{X}$, $\mathsf{T}$
      & all kinds \\
    No longer sold (end of sale)
      & $\mathsf{X}$, $\mathsf{T}$; $\mathsf{I}$ only for recorded customer
        commitments; $\mathsf{L}$ only for supplier-driven changes
      & firmware and software; hardware and process only for
        supplier-driven changes \\
    No longer produced (end of production)
      & $\mathsf{X}$, $\mathsf{T}$
      & firmware and software, preferably on a maintenance stream \\
    Security only
      & $\mathsf{X}$ with the security attribute
      & security patches \\
    Retired (end of service) & none & none \\
    \bottomrule
  \end{tabularx}
\end{table}

\begin{definition}[Gate, publication]
  \label{def:gate}
  Every planned baseline $g$ acts as a \emph{gate}. Its \emph{targeted
  set} at time $t$ is
  \[
    \mathcal{B}_t(g) = \{\, b \in \mathcal{B} : \sigma(b,t) =
    \textsf{targeted} \ \wedge\ g(b) = g \,\},
  \]
  written $\mathcal{B}(g)$ when the time is clear from context.
  \emph{Publishing} $g$ at time $t$ is an atomic operation that sets
  $s(g) = \textsf{released}$ and $t(g) = t$ and merges every branch of
  $\mathcal{B}_t(g)$ into $g$. Before publication, every member of
  $\mathcal{B}_t(g)$ must be verified (\cref{def:state}); members that are
  not are re-targeted to another planned baseline or discarded first
  (\cref{rule:gate}).
\end{definition}

\begin{definition}[Dependency, propagation, launch bundle]
  \label{def:relations}
  The \emph{dependency} relation
  $\mathrm{dep} \subseteq \mathcal{B} \times (\mathcal{B} \cup
  \mathcal{V})$ states that a branch can be delivered only together with
  or after another branch, or only once a baseline (possibly of another
  line) is released. The \emph{propagation} relation
  $\mathrm{prop} \subseteq \mathcal{B} \times \mathcal{B}$ states that
  $b'$ carries the change of $b$ to another stream or line (a back-port or
  forward-port); then $\tau(b') = \tau(b)$. A \emph{launch bundle} is a set
  $\mathcal{G}$ of planned baselines of distinct lines that are published
  together, in one atomic operation, or not at all.
\end{definition}

\begin{definition}[Deployment overlay]
  \label{def:deployment}
  An HPB model may carry an optional \emph{deployment overlay}: for every
  released baseline $v$ and time $t$, a \emph{rollout status}
  $\mathrm{ro}(v,t) \in \{\textsf{pilot}, \textsf{staged}, \textsf{general}\}$
  with, for staged rollouts, the share of the fleet reached; an
  \emph{installed base} $N(v,t) \in \mathbb{N}$ (units or installations
  running $v$); for each customer $c$ of a finite set $\mathcal{K}$, its
  deployed configuration $d(c,t)$; and, for branches, \emph{commitments}
  $\mathrm{cm}(b) \subseteq \mathcal{K} \times (\mathcal{V} \cup
  \mathbb{T})$, i.e., the customers to which $b$ was promised and the
  baseline or date promised.
\end{definition}

For an issue $b$, let $\mathrm{Aff}(b) = \{f(b)\} \cup \mathrm{aff}(b)$
be the baselines in which it is known to occur. The customers exposed to
$b$ at time $t$ are
\[
  \mathcal{K}(b,t) = \{\, c \in \mathcal{K} :
  d(c,t) \cap \mathrm{Aff}(b) \neq \emptyset \,\}.
\]
Customers who have not upgraded to a baseline that contains a fix, and
commitments whose promised baseline has slipped, are obtained in the same
way.

\begin{definition}[HPB model and graph]
  \label{def:model}
  An \emph{HPB model} is a structure
  \[
    H = (\mathcal{P}, \mathrm{parent}, \mathcal{V}, \ell, k, s, t, \beta,
    \delta, \mathrm{app}, \mathcal{B}, \tau, t_o, f, \Phi, \mathrm{aff},
    \iota, S, \sigma, g, m, \rho, \mathrm{dep}, \mathrm{prop},
    \mathbb{G}, A, \alpha),
  \]
  where $\beta$ also determines the streams of every line, $\mathbb{G}$ is
  the set of launch bundles, $A$ is a set of annotations (comments,
  decisions, evidence and approvals) and
  $\alpha\colon A \to \mathcal{P} \cup \mathcal{V} \cup \mathcal{B}$
  attaches each annotation to a node. The relation $\mathrm{closes}$ is
  derived from $m$, and the deployment overlay, when present, is attached
  to $H$. Its \emph{graph} has the baselines and branches as nodes and
  three kinds of structural edges: \emph{precedence} edges between
  immediately preceding baselines (\cref{def:mainline}), \emph{fork} edges
  $f(b) \to b$ and \emph{merge} edges $b \to m(b)$. Dependencies and
  propagations are drawn as annotation edges.
\end{definition}

\subsection{Well-formedness rules}
\label{sec:model-rules}

An HPB model is \emph{well formed} if it satisfies the following rules at
all times. Each rule is followed by the failure it prevents.

\begin{hpbrule}[Single origin]
  \label{rule:origin}
  Every branch has exactly one fork baseline $f(b)$, which is released,
  satisfies $t(f(b)) \le t_o(b)$, is fixed at intake, and is changed only
  by an audited correction.
\end{hpbrule}
\noindent Without a unique origin, ``relative to which product state?''
has no answer and defects cannot be attributed to versions
(\cref{prop:escaped}). Later confirmations extend $\mathrm{aff}(b)$; they
do not create additional forks.

\begin{hpbrule}[Explicit closure]
  \label{rule:closure}
  Every closed branch has exactly one terminal: either a merge into exactly
  one baseline $m(b)$, or a discard with a non-empty rationale $\rho(b)$.
  No branch is closed silently.
\end{hpbrule}
\noindent ``Closed'' without a closing version or a reason is the most
common loss of traceability in practice. A change that must reach several
streams or lines is modeled as one branch per target, linked by
$\mathrm{prop}$, so that a missing propagation is detectable
(\cref{sec:operating}).

\begin{hpbrule}[Time moves forward]
  \label{rule:time}
  For every merged branch, $f(b) \prec m(b)$ and $t(m(b)) \ge t_o(b)$.
\end{hpbrule}
\noindent A baseline cannot contain work requested after it was fixed, and
a branch cannot be resolved by a state that precedes its origin. Since
$\prec$ follows streams and the line forest, merges are confined to
successors of the fork: later baselines of the same stream, of streams
forked after it, or of descendant lines. A defect already resolved in an
existing baseline is discarded with rationale \emph{already resolved in
$v$}, so that change logs only contain work actually incorporated.

\begin{hpbrule}[Gate integrity]
  \label{rule:gate}
  A planned baseline $g$ may be published at time $t$ only if every
  $b \in \mathcal{B}_t(g)$ is verified at $t$, and every baseline or branch
  on which a member depends is released, respectively merged, at $t$ or is
  published or merged in the same operation. If $g$ belongs to a launch
  bundle $\mathcal{G}$, the condition applies to the union of the targeted
  sets of $\mathcal{G}$, and all baselines of $\mathcal{G}$ are published
  together.
\end{hpbrule}
\noindent Every planned baseline becomes a decision point in the sense of
Stage-Gate~\citep{cooper1990stagegate}, applied to releases rather than
projects: committed work is delivered and verified, explicitly deferred or
explicitly abandoned, never left pending behind a shipped version and
never shipped ahead of what it depends on. The rule has content because
publication is the only way to merge; it is checked against the targeted
set \emph{before} publication.

\begin{hpbrule}[Lineage and life cycle]
  \label{rule:lineage}
  A hardware change that fails the interchangeability test creates a new
  line with an explicit parent link (\cref{def:linefork}); release streams
  of any kind do not. References are never edited in place, and parent
  links are never changed once created. New branches and baselines respect
  the life-cycle phase of their line (\cref{tab:phases}).
\end{hpbrule}
\noindent Editing a reference in place would rewrite the history of units
already in the field; the explicit link preserves the ancestry needed for
maintenance and audit (\cref{prop:lineage}), and the phase table keeps a
retired reference from silently accumulating work.

\begin{hpbrule}[Typed work]
  \label{rule:typed}
  Every branch has exactly one type $\tau(b)$. Sub-tasks are steps of the
  branch they belong to and are not branches.
\end{hpbrule}
\noindent Types determine lanes, colors and metrics; a branch with two
types, or a sub-task counted next to its parent, would be counted twice.
Task hierarchies belong to team boards, not to a product-level view.

\begin{hpbrule}[Traceable discussion]
  \label{rule:discussion}
  Every annotation in $A$ is attached by $\alpha$ to exactly one node of the
  model: a line, a baseline or a branch. Annotations and transitions are
  append-only.
\end{hpbrule}
\noindent Decisions and evidence that float in chat threads or meeting
notes cannot be audited; anchoring them to the node they concern puts the
rationale where the diagram is read, and append-only keeps the record
trustworthy.

\begin{remark}
  \label{rem:dag}
  The graph of a well-formed model is acyclic. Precedence edges follow
  $\prec$, which is a forest order; a fork edge followed by a merge edge
  goes from $f(b)$ to $m(b)$ with $f(b) \prec m(b)$ by \cref{rule:time};
  and no edge enters a branch except its fork edge, since sub-tasks are
  steps (\cref{rule:typed}). HPB thus keeps the structural property that
  makes version-control histories readable: every path runs forward in
  precedence. The relation $\mathrm{dep}$ must itself be acyclic, since a
  cycle of dependencies could never be released under \cref{rule:gate}.
\end{remark}

\subsection{Derived properties}
\label{sec:model-properties}

The rules are chosen so that release notes, quality attribution, flow
figures and audit answers are \emph{functions} of the graph rather than
separate records maintained by hand.

\begin{property}[Change log]
  \label{prop:changelog}
  For every baseline $v$, the set
  $C(v) = \{\, b \in \mathcal{B} : \mathrm{closes}(b, v) \,\}$ is the
  content of $v$. The non-empty sets among $\{C(v)\}_{v \in \mathcal{V}}$
  partition the set of merged branches, and $C(v)$ does not change after
  $v$ is released except through audited reopen transitions.
\end{property}
\begin{proof}[Proof sketch]
  By \cref{rule:closure}, each merged branch has exactly one merge baseline,
  hence belongs to exactly one $C(v)$. Merges happen only at publication
  (\cref{def:merge,def:gate}), so no ordinary transition adds a branch to
  $C(v)$ once $v$ is released. Computing all change logs is a single pass
  over $\mathcal{B}$; grouping $C(v)$ by $\tau$ yields release notes
  without manual compilation.
\end{proof}

\begin{property}[Escaped defects]
  \label{prop:escaped}
  For every baseline $v$, the set
  $E(v) = \{\, b \in \mathcal{B} : \tau(b) = \mathsf{X} \ \wedge\ f(b) = v
  \,\}$ contains the issues observed in $v$. The non-empty sets among
  $\{E(v)\}_{v \in \mathcal{V}}$ partition the set of issues, and every
  merged $b \in E(v)$ is resolved in a successor of $v$, i.e.,
  $v \prec m(b)$.
\end{property}
\begin{proof}[Proof sketch]
  By \cref{rule:origin,rule:typed}, each issue has exactly one type and
  exactly one fork baseline, hence belongs to exactly one $E(v)$. The
  second claim is \cref{rule:time} with $f(b) = v$.
\end{proof}
\noindent $E(v)$ attributes \emph{observed} quality to versions; the
baseline $\iota(b)$ in which a defect was \emph{introduced} may be earlier
and requires root-cause analysis. The property holds only if $f(b)$ is the
observed baseline, which is why \cref{rule:origin} fixes it at intake and
why an implementation must store it explicitly (\cref{sec:impl-datamodel};
HPB Markup stores it through \texttt{from}).

\begin{property}[Work in progress]
  \label{prop:wip}
  For every time $t$, the open set
  $W(t) = \{\, b \in \mathcal{B} : \sigma(b,t) \in \{\textsf{open},
  \textsf{targeted}\} \,\}$ is well defined and partitions by type into
  $W_{\mathsf{I}}(t), W_{\mathsf{X}}(t), W_{\mathsf{T}}(t),
  W_{\mathsf{L}}(t)$. Count every maximal interval during which a branch is
  open or targeted as one \emph{visit}. If, over a long observation window,
  the arrival rate $\lambda$ of visits and their mean duration
  $\overline{LT}$ have finite limits and the number of open branches does
  not grow without bound, then the time-average of $|W(t)|$ equals
  $\lambda \cdot \overline{LT}$.
\end{property}
\begin{proof}[Proof sketch]
  $\sigma$ is a total function on the branches opened by time $t$
  (\cref{def:state}), so $W(t)$ is determined, and \cref{rule:typed} makes
  the type partition exact. With visits as customers of a queueing system,
  $|W(t)|$ is the number of customers in the system and the visit length is
  the sojourn time. Little's law~\citep{little1961} then applies under the
  stated limit conditions; it requires no assumption on the distribution of
  arrivals or on the order of service. Merges and discards both count as
  departures, so discarded work must be included when visit lengths are
  averaged. (The arrival rate $\lambda$ here is unrelated to the lane map
  of \cref{eq:lanemap}.)
\end{proof}

\begin{property}[Lineage]
  \label{prop:lineage}
  For every line $P$, the ancestry $P, \mathrm{parent}(P),
  \mathrm{parent}^2(P), \dots$ is finite and is the unique path from a root
  of the line forest to $P$. For every baseline $v$, its predecessors
  $\{u : u \preceq v\}$ form a chain, so that the \emph{cumulative content}
  $\widehat{C}(v) = \bigcup_{u \preceq v} C(u)$ is well defined and follows
  the ancestry of $v$ through streams and lines.
\end{property}
\begin{proof}[Proof sketch]
  By \cref{rule:lineage}, each line has at most one parent, assigned when
  the line is created and never changed afterwards; since a parent must
  exist before its child, ancestry strictly decreases in creation order,
  cannot cycle, and terminates at a root of the forest. For the second
  claim, every baseline has at most one immediate predecessor
  (\cref{def:mainline}), and $\prec$ is well founded because each
  immediate predecessor was created earlier; following immediate
  predecessors from $v$ therefore enumerates $\{u : u \prec v\}$ as a
  single finite chain.
\end{proof}

\noindent As a corollary, a fix merged into a baseline of a maintenance
stream, or into a later baseline of a parent line after a line fork, is
\emph{not} in $\widehat{C}(w)$ for any baseline $w$ of the main stream or
of the child line that does not descend from it. By \cref{rule:closure},
the decision to carry it over is recorded as a second branch linked by
$\mathrm{prop}$, or as an explicit decision on the first one
(\cref{fig:backport}).

\subsection{Ordinal layout}
\label{app:formal-ordinal}

Let $\mathcal{E}_{\le}$ be the finite set of \emph{historical events}
drawn in a diagram, i.e., those with a time $t(e) \le t_{\mathrm{now}}$:
released baselines, branch starts, steps and terminals, with times taken
from the model. For a discarded branch the terminal time is $t_c(b)$. For
a merged branch it is its verification time $t_r(b)$ (\cref{def:state}),
which does not exceed $t_c(b) = t(m(b))$: the check mark is drawn when the
resolution has been verified, and the rise into the baseline shows the
merge at publication. Let $\mathcal{E}_{>}$ be the planned baselines.

\begin{definition}[Ordinal layout]
  \label{def:ordinal}
  An \emph{ordinal layout} assigns a \emph{column} $c(e)$ to every event.
  The historical events receive the columns $1, \dots, |\mathcal{E}_{\le}|$,
  injectively and such that $t(e) < t(e')$ implies $c(e) < c(e')$; events
  with equal times are ordered by a fixed tie-breaking rule (terminals
  before the baseline that closes them, baselines before all other branch
  events, then by identifier). The NOW marker receives the next column,
  $c_{\mathrm{now}} = |\mathcal{E}_{\le}| + 1$, and the planned baselines
  of $\mathcal{E}_{>}$ follow in order of target date. Spacer columns that
  carry no event may be inserted (rule~\ref{lay:tags}). Column $j$ is drawn
  at horizontal position $x = j \cdot w$ for a constant column width $w$.
  Calendar time is shown only as tick labels on a time axis below the
  diagram, at the column of the first event of a calendar period, written
  ``Mon YYYY'' at the first tick and after every change of year.
\end{definition}

A \emph{linear} layout instead uses $x(e) = (t(e) - t_{\min}) / \Delta$
for a fixed time scale $\Delta$. It preserves durations but requires
collision handling for tags and labels and cannot guarantee a minimum
spacing between events (\cref{sec:notation-time}).

% =====================================================================
\subsection{Metric definitions}
\label{app:metrics}

This subsection defines the metrics of \cref{tab:metrics}. $\mathcal{B}$
is the set of branches of a product line (or of several lines, when
stated) and $\mathcal{V}$ its set of baselines. Closed branches split into
merged, $\mathcal{B}^{\mathrm{m}}$, and discarded,
$\mathcal{B}^{\mathrm{d}}$; $\mathcal{B}_\theta$ is the set of branches of
type $\theta$. The closing time $t_c(b)$ is the publication time of $m(b)$
for a merged branch and the time of the discard otherwise
(\cref{def:state}). Windows are intervals $(t_1,t_2]$ of length
$\Delta=t_2-t_1$, and a subscript $(t_1,t_2]$ keeps the branches with
$t_c(b)$ in the window. Lead times are measured per \emph{visit} (each
period during which a branch is open or targeted).

\paragraph{Flow.}
\begin{gather}
  \mathrm{LT}(b) = t_c(b) - t_o(b), \quad b \in \mathcal{B}^{\mathrm{m}} \cup \mathcal{B}^{\mathrm{d}};
  \label{eq:leadtime}\\
  \mathrm{Age}(b,t) = t - t_o(b), \quad b \in W(t);
  \label{eq:age}\\
  W(t) = \{\, b \in \mathcal{B} : \sigma(b,t) \in \{\textsf{open}, \textsf{targeted}\} \,\},
  \quad W_\theta(t) = W(t) \cap \mathcal{B}_\theta;
  \label{eq:wip}\\
  \mathrm{TH}(t_1,t_2) = \bigl|\{\, b \in \mathcal{B}^{\mathrm{m}} : t_1 < t_c(b) \le t_2 \,\}\bigr| \,/\, \Delta .
  \label{eq:throughput}
\end{gather}
Little's law \citep{little1961} states that, in the long run and provided
the averages exist,
\begin{equation}
  \overline{|W|} = \lambda \cdot \overline{\mathrm{LT}},
  \label{eq:little}
\end{equation}
where $\overline{|W|}$ is the time average of $|W(t)|$, $\lambda$ the rate
at which visits start (equal, in a stable system, to the rate at which
they end) and $\overline{\mathrm{LT}}$ the mean lead time per visit. It
makes no assumption about the distribution of arrivals, but it needs a
window long enough and no growing population of branches that never close
(the \emph{eternal branches} of \cref{subsec:method-antipatterns}).

\paragraph{Closure.}
\begin{gather}
  \mathrm{MR} = \frac{|\mathcal{B}^{\mathrm{m}}_{(t_1,t_2]}|}
    {|\mathcal{B}^{\mathrm{m}}_{(t_1,t_2]}| + |\mathcal{B}^{\mathrm{d}}_{(t_1,t_2]}|},
  \qquad \mathrm{DR} = 1-\mathrm{MR};
  \label{eq:mergeratio}\\
  D(b) = \bigl|\{\, v \in \mathcal{V} : s(v)=\textsf{released},\ f(b) \prec v \preceq m(b) \,\}\bigr| ,
  \label{eq:mergedistance}
\end{gather}
where $\prec$ is the precedence of \cref{def:mainline}, crossing the
lineage connector when a branch merges into a child line. Discards
\emph{resolved without a product baseline} are reported apart from
abandoned work.

\paragraph{Releases and quality.}
\begin{gather}
  |C(v)| = \bigl|\{\, b : \mathrm{closes}(b,v) \,\}\bigr| ;
  \label{eq:batch}\\
  E(v) = \{\, b \in \mathcal{B}_{\mathsf{X}} : f(b) = v \,\}, \qquad
  \epsilon(v) = \frac{|E(v)|}{\max(1,|C(v)|)} ;
  \label{eq:escaped}\\
  \eta(v;T) = \frac{\bigl|\{\, b \in E(v) : t_o(b) \le t(v) + T \,\}\bigr|}
    {\int_{t(v)}^{t(v)+T} N(v,t)\,\mathrm{d}t} ;
  \label{eq:exposure}\\
  \mathrm{GS}(g) = t(g) - \hat t(g) .
  \label{eq:gateslip}
\end{gather}
$C(v)$ is the release-note list of $v$. $E(v)$ counts issues \emph{seen}
in $v$; when root-cause analysis finds the baseline $\iota(b)$ that
\emph{introduced} a defect, count by $\iota$ as well. $N(v,t)$ is the
installed base (units running $v$ at time $t$), so $\eta$ is the number of
issues seen within a period $T$ after release per unit-time of operation;
without installed-base data, counting the issues opened within a fixed
period after $t(v)$ is a weaker substitute. In the gate slip, $\hat t(g)$
is the target date of $g$ when its content was first committed; it is
reported with the number of branches moved out of $g$ and discarded from
it.

\paragraph{Maintenance.}
\begin{gather}
  \mathrm{ML}(P;t_1,t_2) =
  \frac{\bigl|\{\, b : f(b)\in \mathcal{V}(P),\ t_1 < t_o(b) \le t_2 \,\}\bigr|}
       {\bigl|\{\, b : f(b)\in \mathcal{V}(P)\cup\mathcal{V}(\mathrm{desc}(P)),\ t_1 < t_o(b) \le t_2 \,\}\bigr|} ;
  \label{eq:maintload}\\
  U = \{\, b \in \mathcal{B}^{\mathrm{m}} \cap \mathcal{B}_{\mathsf{X}} :
    \exists\, v \in \mathrm{aff}(b),\ v \npreceq m(b),\
    \nexists\, b' : (b,b') \in \mathrm{prop} \,\} ,
  \label{eq:unprop}
\end{gather}
where $\mathcal{V}(P)$ is the set of baselines of line $P$,
$\mathrm{desc}(P)$ its descendant lines, $\mathrm{aff}(b)$ the baselines
known to be affected and $\mathrm{prop}$ the ``is a port of'' relation
(\cref{def:relations}); branches with a recorded decision not to port are
removed from $U$. The same maintenance ratio works for a maintenance
stream relative to the main stream of its line.

% =====================================================================
\subsection{Layout algorithm}
\label{app:layout}
\label{sec:impl-layout}

Two layout rules fix the geometry of a diagram completely: ordinal columns
(\ref{lay:ordinal}) fix the horizontal position of every event, and ``one
branch per lane at a time'' (\ref{lay:onelane}) constrains the vertical
position of every branch. \Cref{alg:layout} computes both for one product
line. $f(b)$, $m(b)$ and $g(b)$ are the fork, merge and target baselines
of branch $b$, $\tau(b)$ its type, and $c(\cdot)$ the column of an event
(\cref{def:ordinal}).

Each branch occupies a half-open column interval, from its fork baseline
to its end (merge baseline, discard terminal, target baseline or NOW), so a
branch that merges into a baseline and one that forks from it can share a
lane. Assigning lanes is then the classic \emph{interval partitioning}
problem~\citep{kleinberg2006algo}, solved separately for each type. The
greedy solution uses exactly as many lanes as the maximum number of
branches of that type open at the same moment, which is the minimum
possible. Among the free lanes, a branch already closed by NOW takes the
innermost one and a branch still open or targeted the outermost one, so
long-lived work moves outward and later forks reach their inner lanes
without crossing it (\cref{fig:scale}). The HPB Markup renderer~1.0 packs
lanes first-fit per type and does not yet apply this preference.

\begin{algorithm}[tbp]
  \caption{HPB layout of one product line: ordinal columns and per-type
    lane partitioning.}
  \label{alg:layout}
  \small
  \begin{algorithmic}[1]
    \Require line $P$; baselines $\mathcal{V}(P)$; branches $\mathcal{B}$
      of $P$; present time $t_{\mathrm{now}}$
    \Ensure column $c(e)$ of every event; lane $y(b)$ of every branch
    \State $\mathcal{E} \gets$ released baselines, branch starts, steps and
      terminals with time $\le t_{\mathrm{now}}$
    \State sort $\mathcal{E}$ by (time, tie-breaking rank, identifier)
    \State $c(e_j) \gets j$ for the $j$-th event $e_j$ of $\mathcal{E}$;
      \quad $c_{\mathrm{now}} \gets |\mathcal{E}| + 1$
    \State assign columns $c_{\mathrm{now}}+1, c_{\mathrm{now}}+2, \dots$ to
      the planned baselines in order of target date
    \ForAll{$b \in \mathcal{B}$} \Comment{lane interval $[s_b, e_b)$}
      \State $s_b \gets c(f(b))$
      \State $e_b \gets c(m(b))$ if merged; $c(\text{terminal})$ if
        discarded; $c(g(b))$ if targeted; $c_{\mathrm{now}}$ otherwise
    \EndFor
    \ForAll{$\theta \in \{\mathsf{I}, \mathsf{L}, \mathsf{X}, \mathsf{T}\}$}
      \State $Q \gets$ branches of type $\theta$ sorted by $(s_b, e_b)$
      \State $k_\theta \gets$ maximum number of intervals of $Q$ sharing a
        point \Comment{sweep over sorted endpoints}
      \State $F \gets \{1, \dots, k_\theta\}$ \Comment{free lanes,
        ordered set}
      \State $H \gets \emptyset$ \Comment{busy lanes, min-heap on $e_b$}
      \ForAll{$b \in Q$ in order}
        \While{$H \neq \emptyset$ \textbf{and} $\min_{H} e \le s_b$}
          \State move the lane of the minimum of $H$ back to $F$
        \EndWhile
        \If{$b$ is open or targeted at $t_{\mathrm{now}}$}
          \State $\ell_b \gets \max F$ \Comment{outermost free lane}
        \Else
          \State $\ell_b \gets \min F$ \Comment{innermost free lane}
        \EndIf
        \State remove $\ell_b$ from $F$
        \State insert $(e_b, \ell_b)$ into $H$
      \EndFor
    \EndFor
    \ForAll{$b \in \mathcal{B}$} \Comment{stack lanes outward from the
      main line}
      \State $y(b) \gets
        \begin{cases}
          +\ell_b & \tau(b) = \mathsf{I}\\
          +(k_{\mathsf{I}} + \ell_b) & \tau(b) = \mathsf{L}\\
          -\ell_b & \tau(b) = \mathsf{X}\\
          -(k_{\mathsf{X}} + \ell_b) & \tau(b) = \mathsf{T}
        \end{cases}$
    \EndFor
  \end{algorithmic}
\end{algorithm}

With $n$ events and branches, the layout runs in $O(n \log n)$ time and
$O(n)$ space; drawing glyphs, tags and labels afterwards is linear. Every
sort ends with the identifier, so the same model at the same present time
always gives the same diagram. Two variants keep long histories readable:
\emph{collapsed history} hides the merged branches of all but the most
recent released baselines and replaces them by count badges
(\cref{fig:scale}), and \emph{type filters} remove the branches, and the
lanes, of deselected types. An optional \emph{compaction} that lets
consecutive events of one kind within a day, week or month share a column
bends rule~\ref{lay:ordinal}, and a view that uses it must say so.

% =====================================================================
\subsection{Prototype implementation: data model and status}
\label{app:impl}
\label{sec:impl-datamodel}

\Cref{fig:datamodel} shows the overlay schema of the prototype described
in \cref{sec:implementation}. The existing trackers stay as they are.
\texttt{products} holds the product lines, and \texttt{parent\_product\_id}
the family tree of references. \texttt{product\_versions} holds the
baselines with kind, status and date; a hardware baseline that creates a
new reference stores it in \texttt{result\_product\_id}, which is the line
fork of \cref{rule:lineage}. \texttt{product\_plan\_items} attaches one
tracker item, by (\texttt{item\_type}, \texttt{item\_id}), to one line:
the pair gives the branch type, \texttt{version\_id} holds the target
while the branch is open and the merge baseline once it is merged, and
\texttt{outcome}, \texttt{closed\_at} and \texttt{note} record the terminal
and the rationale of a discard. \texttt{product\_plan\_comments} attaches
each comment to exactly one node, so \cref{rule:discussion} is enforced by
the schema.

% F9 -- Data model of the reference implementation: an overlay on existing
% trackers. Requires hpb.sty v1.0 (palette, styles, layers of hpbgraph).
% Units: hpbgraph with col = lane = 1cm, i.e. plain centimetres.
\begin{figure}[t]
\centering
\resizebox{\linewidth}{!}{%
\begin{hpbgraph}[col=1cm, lane=1cm]
  \tikzset{
    dm/head/.style  = {anchor=north west, inner xsep=4pt, inner ysep=3pt,
                       align=left, text=hpbInk,
                       font=\sffamily\fontsize{7}{8.6}\selectfont},
    dm/row/.style   = {anchor=north west, inner xsep=4pt, inner ysep=0pt,
                       minimum height=0.3cm, text height=5pt, text depth=1.6pt,
                       align=left, text=hpbInk,
                       font=\fontsize{6.5}{7.8}\selectfont},
    dm/frame/.style = {draw=hpbInk, line width=1pt, rounded corners=4pt},
    dm/fk/.style    = {draw=hpbInk, line width=0.9pt, rounded corners=5pt,
                       -{Stealth[length=5pt,width=4.5pt]}},
    dm/ref/.style   = {dm/fk, dash pattern=on 4pt off 2.5pt},
    dm/lbl/.style   = {font=\sffamily\fontsize{6.5}{7.8}\selectfont, text=hpbInk,
                       fill=white, inner sep=1.2pt, align=center},
    dm/mult/.style  = {font=\sffamily\bfseries\fontsize{6.5}{7.8}\selectfont,
                       text=hpbMuted, inner sep=1pt},
  }
  % ---- entity macro (local to this figure) --------------------------------
  % \dment{<id>}{<x,y top-left>}{<width cm>}{<fill>}{<table>}{<model element>}
  %       {<key/field, key/field, ...>}   key in {PK, FK, (empty)}
  % Defines nodes <id>-h (header), <id>-1, <id>-2, ... (one per field) and
  % the corners <id>-nw, <id>-se.
  \def\dmPK{PK}%
  \newcommand\dmkey[1]{\edef\dmk{#1}\makebox[0.42cm][l]{\sffamily\bfseries
    \ifx\dmk\dmPK\dmk\else\textcolor{hpbMuted}{\dmk}\fi}}
  \newcommand\dment[7]{%
    \pgfmathsetlengthmacro\dmw{#3cm-8pt}%
    \node[dm/head, text width=\dmw] (#1-h) at (#2)
      {\textbf{\ttfamily\fontsize{7.5}{9}\selectfont #5}\\[-0.5pt]{\itshape #6}};
    \foreach \k/\f [count=\i] in {#7}{%
      \node[dm/row, text width=\dmw] (#1-\i)
        at ([yshift={-0.08cm-(\i-1)*0.3cm}]#1-h.south west)
        {\dmkey{\k}\texttt{\f}};
      \xdef\dmlast{\i}}%
    \coordinate (#1-se) at ([yshift=-0.1cm]#1-\dmlast.south west -| #1-h.east);
    \coordinate (#1-nw) at (#1-h.north west);
    \begin{pgfonlayer}{background}
      \begin{scope}
        \clip[rounded corners=4pt] (#1-nw) rectangle (#1-se);
        \fill[white] (#1-nw) rectangle (#1-se);
        \fill[#4] (#1-nw) rectangle (#1-h.south east);
      \end{scope}
    \end{pgfonlayer}
    \draw[draw=hpbInk, line width=0.6pt] (#1-h.south west) -- (#1-h.south east);
    \draw[dm/frame] (#1-nw) rectangle (#1-se);}

  % ---- HPB tables -------------------------------------------------------------
  \dment{pl}{0,0}{3.7}{hpbMain}{products}{product line $P$}
    {PK/id, /model, /manufacturer\_ref, FK/parent\_product\_id}
  \dment{bl}{0,-4.25}{3.7}{hpbMain}{product\_versions}{baseline $v$}
    {PK/id, FK/product\_id, /name, /kind, /status,
     /target\_date, /released\_at, FK/result\_product\_id}
  \dment{pi}{5.9,-0.45}{3.9}{hpbHw!45}{product\_plan\_items}
    {state of branch $b$ on line $P$}
    {PK/id, FK/product\_id, /item\_type, /item\_id, FK/version\_id,
     /outcome, /closed\_at, /note}
  \dment{cm}{5.9,-4.95}{3.9}{hpbDisc}{product\_plan\_comments}
    {discussion on a node (R7)}
    {PK/id, FK/product\_id, /target\_type, /target\_id, /body,
     {/created\_by, created\_at}}

  % ---- existing trackers (referenced, not migrated) -------------------------
  \draw[draw=hpbMuted, line width=0.8pt, rounded corners=5pt,
        dash pattern=on 3pt off 2pt, fill=hpbGuide!25]
    (11.75,-0.1) rectangle (15.95,-5.55);
  \node[font=\sffamily\bfseries\fontsize{7}{8.4}\selectfont, text=hpbInk,
        anchor=north west] at (11.85,-0.15) {existing trackers};
  \node[font=\sffamily\itshape\fontsize{6.5}{7.8}\selectfont, text=hpbMuted,
        anchor=north west, align=left] at (11.85,-0.5)
    {branch $b$ with type $\tau(b)$};
  \foreach \t/\c/\l/\y in {improvements/hpbImp/I/-1.05,
                           issues/hpbIss/X/-2.15,
                           tasks/hpbTask/T/-3.25,
                           lab\_items/hpbLab/L/-4.35}{
    \begin{scope}
      \clip[rounded corners=4pt] (12.0,\y) rectangle (15.7,\y-0.9);
      \fill[white] (12.0,\y) rectangle (15.7,\y-0.9);
      \fill[\c] (12.0,\y) rectangle (15.7,\y-0.42);
    \end{scope}
    \draw[dm/frame] (12.0,\y) rectangle (15.7,\y-0.9);
    \draw[draw=hpbInk, line width=0.6pt] (12.0,\y-0.42) -- (15.7,\y-0.42);
    \node[hpb/event, fill=\c, minimum size=9.5pt, line width=1.1pt,
          font=\sffamily\bfseries\fontsize{6.5}{7.8}\selectfont]
      at (12.3,\y-0.21) {\l};
    \node[font=\ttfamily\bfseries\fontsize{7.5}{9}\selectfont, text=hpbInk,
          anchor=west, inner sep=0pt] at (12.62,\y-0.21) {\t};
    \node[font=\ttfamily\fontsize{6.5}{7.8}\selectfont, text=hpbInk,
          anchor=west, inner sep=0pt] at (12.12,\y-0.66)
      {id, reference, title, status};
  }
  \coordinate (tr-w)  at (11.75,-2.15);
  \coordinate (tr-s)  at (13.85,-5.55);

  % ---- foreign keys (solid): referencing -> referenced table ---------------
  % products.parent_product_id -> products.id   (line forest)
  \draw[dm/fk] (pl-4.west) -- ++(-0.35,0) |- (pl-1.west);
  \node[dm/lbl, rotate=90] at ([xshift=-0.35cm]pl-1.west |- 0,-1.42) {parent};
  % product_versions.result_product_id -> products   (line fork: new reference)
  \draw[dm/fk] (bl-8.west) -- ++(-0.75,0) |- (pl-h.west);
  \node[dm/lbl, rotate=90] at (-0.75,-4.6) {line fork (new reference)};
  \node[dm/mult, anchor=north east] at ([xshift=-1pt]bl-8.west) {*};
  \node[dm/mult, anchor=south east] at ([xshift=-2pt,yshift=2.5pt]pl-h.west) {0..1};
  % product_versions.product_id -> products   (baselines of a line)
  \draw[dm/fk] (1.85,-4.25) -- (1.85,0 |- pl-se);
  \node[dm/lbl, anchor=west] at (1.9,-3.15) {product\_id};
  \node[dm/mult, anchor=north west] at ([xshift=3pt,yshift=-1pt]1.85,0 |- pl-se) {1};
  \node[dm/mult, anchor=south west] at ([xshift=1pt]1.85,-4.25) {*};
  % product_plan_items.product_id -> products
  \draw[dm/fk] (pi-2.west) -- (4.75,0 |- pi-2.west) |- (pl-2.east);
  \node[dm/mult, anchor=south east] at ([xshift=-1pt]pi-2.west) {*};
  \node[dm/mult, anchor=south west] at ([xshift=2pt,yshift=2.5pt]pl-2.east) {1};
  % product_plan_items.version_id -> product_versions   (g(b) or m(b))
  \draw[dm/fk] (pi-5.west) -- (4.75,0 |- pi-5.west) |- (bl-3.east);
  \node[dm/mult, anchor=south east] at ([xshift=-1pt]pi-5.west) {*};
  \node[dm/mult, anchor=south west] at ([xshift=2pt,yshift=2.5pt]bl-3.east) {0..1};
  \node[dm/lbl, anchor=west, align=left] at (4.8,-4.4)
    {target $g(b)$ or\\closing $m(b)$};

  % ---- polymorphic references (dashed): (type, id) pairs ---------------------
  \draw[dm/ref] (pi-4.east) -- (tr-w |- pi-4.east);
  \node[dm/lbl, anchor=south] at ([yshift=1pt]10.78,0 |- pi-4.east)
    {item\_type,\\item\_id};
  \node[dm/lbl, anchor=north, font=\sffamily\itshape\fontsize{6.5}{7.8}\selectfont,
        text=hpbMuted] at ([yshift=-3pt]10.78,0 |- pi-4.east)
    {referenced,\\not migrated};
  % comments -> branch item, baseline (or product)
  \draw[dm/ref] (cm-4.east) -| (tr-s);
  \draw[dm/ref] (cm-nw |- bl-7.east) -- (bl-7.east);
  \node[dm/lbl, anchor=north west, align=left] at ([yshift=-2pt]cm-se -| cm-nw)
    {target\_type $\in$ \{product, version, improvement,\\
     issue, task, lab\}, identified by target\_id};
\end{hpbgraph}}
\caption{Data model of the prototype implementation. HPB is an
overlay: improvements, issues, tasks and lab items stay in their trackers and
are referenced, not migrated. Solid arrows are foreign keys; dashed arrows are
references by (type, identifier).}
\label{fig:datamodel}
\end{figure}


Some attributes are \emph{derived} today instead of stored: the opening
time of a branch is the creation time of its tracker item, and the fork
baseline is the latest release at that time. That default is right for
improvements, tasks and lab work, but wrong for issues, whose fork
baseline is the version in which the defect was observed
(\cref{prop:escaped}). The design therefore adds a stored
\texttt{fork\_version\_id}, mandatory for issues and with an
\emph{inferred} flag, and an append-only log of transitions with actor,
approver and time, which makes audited reopens possible and keeps the
original target date from which gate slip is measured. Streams,
dependencies, propagations, launch bundles, life-cycle milestones and the
deployment overlay need further tables.

\Cref{tab:impl-status} separates what the design specifies from what the
prototype does today. The prototype still accepts a discard without a
rationale, offers released baselines when an item is closed as merged,
merges the whole targeted set without checking verification, derives the
fork baseline instead of storing it, and packs lanes per side of the main
line instead of per type. It therefore does not yet enforce
\cref{rule:closure,rule:time,rule:gate}; until it does, those violations
are at least visible in the diagram.

\begin{table}[tbp]
  \centering
  \small
  \caption{Prototype in a product management portal: designed versus
    built. \checkmark: built; $\sim$: partly built; --: designed only.}
  \label{tab:impl-status}
  \begin{tabularx}{\linewidth}{@{}>{\raggedright\arraybackslash}X
      >{\centering\arraybackslash}p{1.4cm}
      >{\raggedright\arraybackslash}p{6.2cm}@{}}
    \toprule
    \textbf{Element of the design} & \textbf{Status} & \textbf{Remark} \\
    \midrule
    Overlay tables, references to tracker items & \checkmark & no migration of trackers \\
    Derived mode & \checkmark & read-only reconstruction \\
    Line family tree and line forks & \checkmark & \texttt{parent\_product\_id}, \texttt{result\_product\_id} \\
    Stored observed fork baseline, inferred flag & -- & needed for escaped defects per baseline \\
    Release streams, configurations, applies-to & -- & \\
    Dependencies, propagations, launch bundles & -- & \\
    Life-cycle milestones, deployment overlay & -- & \\
    Append-only transition log with approvers & -- & state tables only \\
    R1 single origin & $\sim$ & by derivation, not at intake \\
    R2 explicit closure & $\sim$ & rationale not yet mandatory \\
    R3 time moves forward & $\sim$ & released baselines still offered \\
    R4 gate integrity & $\sim$ & batch merge without verification \\
    R5 lineage and life cycle & \checkmark & life-cycle phases not checked \\
    R6 typed work & \checkmark & type from the source tracker \\
    R7 traceable discussion & $\sim$ & comments on nodes; not append-only \\
    Ordinal columns & \checkmark & \\
    Lane partitioning per type (\cref{alg:layout}) & -- & today: per side of the main line \\
    \bottomrule
  \end{tabularx}
\end{table}

% flush this appendix's floats before the next appendix
\clearpage
